\section{Symmetric Laurent polynomials}\label{macdonaldsection}

A canonical reference for symmetric polynomials is~\cite{M2}. We choose to give a brief overview of the tools that we need from~\cite{M2}, in order to f\/ix terminology and to introduce lesser known objects, such as signatures and Macdonald Laurent polynomials.

\subsection{Partitions, signatures and symmetric Laurent polynomials}

A \textit{partition} is a f\/inite sequence of weakly decreasing nonnegative integers $\lambda = (\lambda_1 \geq \lambda_2 \geq \cdots\geq \lambda_k)$, $\lambda_i\in\Z_{\geq 0}$ $\forall\, i$. We identify partitions that dif\/fer by trailing zeroes; for example, $(4, 2, 2, 0, 0)$ and $(4, 2, 2)$ are the same partition. We def\/ine the \textit{size} of $\lambda$ to be the sum $|\lambda| \myeq \lambda_1 + \dots + \lambda_k$, and its \textit{length} $\ell(\lambda)$ to be the number of strictly positive elements of it.
The \textit{dominance order} for partitions is a partial order given by letting $\mu \leq \lambda$ if $|\mu| = |\lambda|$ and $\mu_1 + \cdots + \mu_i \leq \lambda_1 + \cdots + \lambda_i$ for all $i$. As usual, we let $\mu < \lambda$ if $\mu \leq \lambda$ and $\mu \neq \lambda$.

Partitions can be graphically represented by their \textit{Young diagrams}. The Young diagram of partition $\lambda$ is the array of boxes with coordinates $(i, j)$ with $1\leq j\leq \lambda_i$, $1\leq i\leq\ell(\lambda)$, where the coordinates are in matrix notation (row labels increase from top to bottom and column labels increase from left to right), see Fig.~\ref{youngdiagram1}.

\begin{figure}[h!]\centering
\ytableausetup{centertableaux}
\begin{ytableau}
\ & & & & \\
 & & & \\
 & & s & \\
&
\end{ytableau}
\caption{Young diagram for the partition $\lambda = (5, 4, 4, 2)$. Square $s = (3, 3)$ has arm length, arm colength, leg length and leg colength given by $a(s) = 1$, $a'(s) = 2$, $l(s) = 0$, $l'(s) = 2$.}\label{youngdiagram1}
\end{figure}

A \textit{signature} is a sequence of weakly decreasing integers $\lambda = (\lambda_1 \geq \lambda_2 \geq \dots\geq \lambda_k)$, $\lambda_i\in\Z$ $\forall\, i$. A \textit{positive signature} is a signature whose elements are all nonnegative. The \textit{length} of a~signature, or positive signature, is the number~$k$ of elements of it. Positive signatures which dif\/fer by trailing zeroes are not identif\/ied, in contrast to partitions; for example, $(4, 2, 2, 0, 0)$ and $(4, 2, 2)$ are dif\/ferent positive signatures, the f\/irst of length $5$ and the second of length $3$. We shall denote~$\GT_N$ (resp.~$\GTp_N$) the set of signatures (resp. positive signatures) of length~$N$. Evidently~$\GTp_N$ can be identif\/ied with the set of all partitions of length $\leq N$. Under this identif\/ication, we are allowed to talk about the Young diagram of a positive signature $\lambda\in\GTp_N$, its size, the dominance order, and other attributes that are typically associated to partitions. Note, however, that length is def\/ined dif\/ferently for partitions and for positive signatures.

Let us now switch to notions pertaining to symmetric (Laurent) polynomials. Fix a positive integer $N$. Consider the f\/ield $F = \C(q, t)$ and recall the algebra $\Lambda_F[x_1, \dots, x_N]$ of symmetric polynomials on the variables $x_1, \dots, x_N$ with coef\/f\/icients in $F$. For any $m\in\Z_{\geq 0}$, recall also the subalgebra $\Lambda_F^m[x_1, \dots, x_N]$ of symmetric polynomials on $x_1, \dots, x_N$ that are homogeneous of degree $m$; then
\begin{gather*}
\Lambda_F[x_1, \dots, x_N] = \bigoplus_{m\geq 0}{\Lambda_F^m[x_1, \dots, x_N]}.
\end{gather*}
We also denote by $\Lambda_F[x_1^{\pm}, \dots, x_N^{\pm}]$ the algebra of symmetric (with respect to the transpositions $x_i \leftrightarrow x_{i+1}$ for $i = 1, \dots, N-1$) Laurent polynomials in the variables $x_1, \dots, x_N$.

The connection between partitions/signatures and symmetric polynomials comes from the observation that $\dim_F \left(\Lambda_F^m[x_1, \dots, x_N]\right)$ is the number of partitions of size $m$ and length $\leq N$, or equivalently the number of positive signatures of size $m$ and length $N$. A basis for the space $\Lambda_F^m[x_1, \dots, x_N]$ is given by the \textit{monomial symmetric polynomials} $m_{\lambda}(x_1, \dots, x_N)$, with $|\lambda| = m$, $\ell(\lambda)\leq N$, def\/ined by
\begin{gather*}
m_{\lambda}(x_1, \dots, x_N) \myeq \sum_{\mu\in S_N\cdot\lambda}{x_1^{\mu_1}\cdots x_N^{\mu_N}},
\end{gather*}
where $S_N\cdot\lambda$ is the orbit of $\lambda$ under the permutation action of $S_N$, and the sum runs over distinct elements $\mu$ of that orbit. It is implied that $\{m_{\lambda}(x_1, \dots, x_N)\colon \ell(\lambda) \leq N\}$ is a basis of $\Lambda_F[x_1, \dots, x_N]$.

\subsection{Macdonald polynomials and Macdonald characters}\label{macdonaldpolysection}

\begin{dfprop}[{\cite[Chapter~VI, Sections~3, 4, 9]{M2}}]\label{def:macdonaldpolys}
The \textit{Macdonald polynomials} $P_{\lambda}(x_1, \dots, x_N; q, t)$, for partitions $\lambda$ with $\ell(\lambda) \leq N$, are the unique elements of $\Lambda_F[x_1, \dots, x_N]$ satisfying the following two properties
\begin{itemize}\itemsep=0pt
	\item \textit{Triangular decomposition}: $\displaystyle P_{\lambda}(x_1, \dots, x_N; q, t) = m_{\lambda} + \sum\limits_{\mu\colon \mu < \lambda}{c_{\lambda, \mu}m_{\mu}}$, for some $c_{\lambda, \mu}\in F$, and the sum is over partitions $\mu$ with $\ell(\mu) \leq N$, and $\mu < \lambda$ in the dominance order.

	\item \textit{Orthogonality relation}: Let $[\cdot]_0 \colon F[x_1^{\pm}, \dots, x_N^{\pm}] \rightarrow F$ be the constant term map
\begin{gather*}
\left[ \sum_{\lambda = (\lambda_1 \geq \cdots \geq \lambda_N) \in \Z^N}{a_{\lambda}x_1^{\lambda_1}\cdots x_N^{\lambda_N}} \right]_0 = a_{(0, \dots, 0)}.
\end{gather*}
The Macdonald polynomials are orthogonal with respect to the inner product $(\cdot, \cdot)_{q, t}$ on $\Lambda_F[x_1^{\pm}, \dots, x_N^{\pm}]$ given by $\displaystyle (f, g)_{q, t} \myeq [f(x_1, \dots, x_N)g\big(x_1^{-1}, \dots, x_N^{-1}\big)\Delta_{q, t}]_0$, where
\[
\Delta_{q, t} \myeq \prod_{1\leq i\neq j\leq N}\prod_{k=0}^{\infty}{\frac{1 - q^{k}x_ix_j^{-1}}{1 - q^{k}tx_ix_j^{-1}}}.
\]
\end{itemize}
\end{dfprop}

Note that $P_{\varnothing}(q, t) = 1$. If $N < \ell(\lambda)$, we set $P_{\lambda}(x_1, \dots, x_N; q, t) \myeq 0$ for convenience.

When we are talking about Macdonald polynomials and some of their properties which hold regardless of the number $N$ of variables, as long as $N$ is large enough, we simply write $P_{\lambda}(q, t)$ instead of $P_{\lambda}(x_1, \dots, x_N; q, t)$.

From the triangular decomposition of Macdonald polynomials, $P_{\lambda}(q, t)$ is a homogeneous polynomial of degree $|\lambda|$. Moreover $\{P_{\lambda}(x_1, \dots, x_N) \colon \ell(\lambda) \leq N\}$ is a basis of $\Lambda_F[x_1, \dots, x_N]$. Finally, we have the following \textit{index stability} property:
\begin{gather}\label{macdonaldsignatures1}
P_{(\lambda_1 + 1, \dots, \lambda_N + 1)}(x_1, \dots, x_N; q, t) = (x_1\cdots x_N)\cdot P_{\lambda}(x_1, \dots, x_N; q, t).
\end{gather}

As pointed out before, the set of partitions of length $\leq N$ is in bijection with $\GTp_N$. Thus we can index the Macdonald polynomials by positive signatures rather than by partitions: for any $\lambda\in\GTp_N$, we let $P_{\lambda}(x_1, \dots, x_N; q, t)$ be the Macdonald polynomial corresponding to the partition associated to $\lambda$. We can slightly extend the def\/inition above and introduce Macdonald Laurent polynomials $P_{\lambda}(x_1, \dots, x_N; q, t)$ for any $\lambda\in\GT_N$. Let $\lambda\in\GT_N$ be arbitrary. If $\lambda_N \geq 0$, then $\lambda\in\GTp_N$ and $P_{\lambda}(x_1, \dots, x_N; q, t)$ is already def\/ined. If $\lambda_N < 0$, choose $m\in\N$ such that $\lambda_N + m \geq 0$ and so $(\lambda_1 + m, \dots, \lambda_N + m)\in\GTp_N$. Then def\/ine
\begin{gather}\label{macdonaldsignatures11}
P_{\lambda}(x_1, \dots, x_N; q, t) \myeq (x_1\cdots x_N)^{-m}\cdot P_{(\lambda_1 + m, \dots, \lambda_N + m)}(x_1, \dots, x_N; q, t).
\end{gather}
By virtue of the index stability property, the Macdonald (Laurent) polynomial $P_{\lambda}(x_1, \dots, x_N; q, t)$ is well-def\/ined and does not depend on the value of $m$ that we choose. For simplicity, we call $P_{\lambda}(x_1, \dots, x_N; q, t)$ a Macdonald polynomial, whether $\lambda\in\GTp_N$ or not. In a similar fashion, we can def\/ine monomial symmetric polynomials $m_{\lambda}$, for any $\lambda\in\GT_N$.

Recall the def\/initions of the \textit{arm-length, arm-colength, leg-length, leg-colength} $a(s)$, $a'(s)$, $l(s)$, $l'(s)$ of the square $s = (i, j)$ of the Young diagram of $\lambda$, given by $a(s) = \lambda_i - j$, $a'(s) = j - 1$, $l(s) = \lambda_j' - i$, $l'(s) = i - 1$;
we note that $\lambda_j' = |\{ i\colon \lambda_i \geq j \} |$ is the length of the $j$th part of the \textit{conjugate partition} $\lambda'$, see Fig.~\ref{youngdiagram1}.

We use terminology from $q$-analysis, see Appendix~\ref{app:qtheory}; particularly we use the def\/inition of $q$-Pochhammer symbols $(z; q)_n \myeq \prod\limits_{i=0}^{n-1}{(1 - zq^i)}$ and $(z; q)_{\infty} \myeq \prod\limits_{i=0}^{\infty}{(1 - zq^i)}$.

For $\lambda\in\bigsqcup\limits_{N\geq 0}{\GTp_N}$, def\/ine the \textit{dual Macdonald polynomials} $Q_{\lambda}(q, t)$ as the following normaliza\-tion of Macdonald polynomials
\begin{gather}\label{PtoQ}
Q_{\lambda}(q, t) \myeq b_{\lambda}(q, t)P_{\lambda}(q, t), \qquad b_{\lambda}(q, t) \myeq \prod_{s\in\lambda}{\frac{1 - q^{a(s)}t^{l(s) + 1}}{1 - q^{a(s) + 1}t^{l(s)}}}.
\end{gather}

The \textit{complete homogeneous symmetric $($Macdonald$)$ polynomials} $g_0 = 1, g_1, g_2, \dots$ are the one-row dual Macdonald polynomials:
\begin{gather}\label{gtoQ}
g_n(q, t) \myeq Q_{(n)}(q, t) = \frac{(q; q)_n}{(t; q)_n}P_{(n)}(q, t).
\end{gather}
For convenience, we also set $g_n(q, t) \myeq 0$, $\forall\, n < 0$.

Now we come to several important theorems on Macdonald polynomials, which will be our main tools.

\begin{thm}[{index-argument symmetry; \cite[Chapter~VI, Property~6.6]{M2}}]\label{symmetry}
Let $N\in\N$, $\lambda, \mu\in\GTp_N$, then
\begin{gather*}
\frac{P_{\lambda}\big(q^{\mu_1}t^{N-1}, q^{\mu_2}t^{N-2}, \dots, q^{\mu_N}; q, t\big)}{P_{\lambda}\big(t^{N-1}, t^{N-2}, \dots, 1; q, t\big)} = \frac{P_{\mu}\big(q^{\lambda_1}t^{N-1}, q^{\lambda_2}t^{N-2}, \dots, q^{\lambda_N}; q, t\big)}{P_{\mu}\big(t^{N-1}, t^{N-2}, \dots, 1; q, t\big)}.
\end{gather*}
\end{thm}

\begin{thm}[{evaluation identity; \cite[Chapter~VI, equations~(6.11) and~(6.11$'$)]{M2}}]\label{evaluation}
Let $N\in\N$, $\lambda\in\GTp_N$, then
\begin{gather*}
P_{\lambda}\big(t^{N-1}, t^{N-2}, \dots, 1; q, t\big) = t^{n(\lambda)}\prod_{1\leq i<j\leq N}{\frac{\big(q^{\lambda_i - \lambda_j}t^{j-i}; q\big)_{\infty}(t^{j-i+1}; q)_{\infty} }{ \big(q^{\lambda_i - \lambda_j}t^{j-i+1}; q\big)_{\infty} (t^{j-i}; q)_{\infty}}}\\
\hphantom{P_{\lambda}\big(t^{N-1}, t^{N-2}, \dots, 1; q, t\big)}{} = t^{n(\lambda)}\prod_{s\in\lambda}{\frac{1 - q^{a'(s)}t^{N - l'(s)}}{1 - q^{a(s)}t^{l(s)+1}}},
\end{gather*}
where and $n(\lambda) \myeq \lambda_2 + 2\lambda_3 + \cdots + (N-1)\lambda_N$.
The first equality holds, more generally, for any signature $\lambda\in\GT_N$ by virtue of the definition~\eqref{macdonaldsignatures11} of Macdonald Laurent polynomials $P_{\lambda}(x_1, \dots, x_N; q, t)$.
\end{thm}

From the second equality of Theorem~\ref{evaluation} and the def\/inition of dual Macdonald polynomials, we obtain

\begin{cor}\label{evaluationcor}
Let $N\in\N$, $\lambda\in\GTp_N$, then
\begin{gather*}
Q_{\lambda}\big(t^{N-1}, t^{N-2}, \dots, 1; q, t\big) = t^{n(\lambda)}\prod_{s\in\lambda}{\frac{1 - q^{a'(s)}t^{N - l'(s)}}{1 - q^{a(s)+1}t^{l(s)}}}.
\end{gather*}
\end{cor}

Since the Macdonald polynomial $P_{\lambda}(x_1, x_2, \dots, x_N; q, t)$ is symmetric in $x_1, x_2, \dots, x_N$, it is also a symmetric polynomial on $x_2, \dots, x_N$; thus it is a linear combination of Macdonald polynomials $P_{\mu}(x_2, \dots, x_N; q, t)$ with coef\/f\/icients in $F[x_1]$. More precisely, we have the so-called \textit{branching rule for Macdonald polynomials}:

\begin{thm}[{branching rule; \cite[Chapter~VI, equation~(7.13$'$), Example~2(b) on p.~342]{M2}}]\label{branchingmacdonald}
Let $N\in\N$, $\lambda\in\GTp_N$, then
\begin{gather*}
P_{\lambda}(x_1, x_2, \dots, x_N; q, t) = \sum_{\mu \in\GTp_{N-1}\colon \mu\prec\lambda}{\psi_{\lambda/\mu}(q, t)x_1^{|\lambda| - |\mu|}P_{\mu}(x_2, \dots, x_N; q, t)},
\end{gather*}
where the branching coefficients are
\begin{gather*}
\psi_{\lambda/\mu}(q, t) \myeq \prod_{1\leq i \leq j\leq N-1}\frac{\big(q^{\mu_i - \mu_j}t^{j-i+1}; q\big)_{\infty}\big(q^{\lambda_i - \lambda_{j+1}}t^{j-i+1}; q\big)_{\infty} }
{\big(q^{\lambda_i - \mu_j}t^{j-i+1}; q\big)_{\infty}\big(q^{\mu_i - \lambda_{j+1}}t^{j-i+1}; q\big)_{\infty} }\\
\hphantom{\psi_{\lambda/\mu}(q, t) \myeq \prod_{1\leq i \leq j\leq N-1}}{} \times \frac{\big(q^{\lambda_i - \mu_j+1}t^{j-i}; q\big)_{\infty}\big(q^{\mu_i - \lambda_{j+1} + 1}t^{j-i}; q\big)_{\infty}}{\big(q^{\mu_i - \mu_j+1}t^{j-i}; q\big)_{\infty}\big(q^{\lambda_i - \lambda_{j+1}+1}t^{j-i}; q\big)_{\infty}} ,
\end{gather*}
and the sum is over positive signatures $\mu\in\GTp_{N-1}$ that satisfy the interlacing constraint
\begin{gather*}
\lambda_N \leq \mu_{N-1} \leq \lambda_{N-1} \leq \cdots \leq \mu_1 \leq \lambda_1,
\end{gather*}
which is written succinctly as $\mu\prec\lambda$.
\end{thm}

Observe that $\psi_{\lambda/\mu}(q, t) > 0$, whenever $q, t \in (0, 1)$. By applying the branching rule several times, we can deduce the following.

\begin{cor}\label{cor:branchingrule}
The coefficients $c_{\lambda, \mu} = c_{\lambda, \mu}(q, t)$ in the expansion
\begin{gather*}
P_{\lambda}(q, t) = \sum_{\mu}{c_{\lambda, \mu}m_{\mu}}
\end{gather*}
are such that $c_{\lambda, \lambda} = 1$ and $c_{\lambda, \mu} \geq 0$, whenever $q, t \in (0, 1)$. Moreover, $c_{\lambda, \mu} = 0$ unless $\lambda \geq \mu$.
\end{cor}

